\begin{abstract}
This study compares feedback from a coverage tool (JaCoCo) and a mutation testing tool (PITest) on how well each helps students write tests that catch real faults. In a crossover experiment, 24 participants each had 20 minutes to add JUnit tests to two IBM CodeNet programs, using one tool per program. The outcome measure was the gain in detection rate on held-out Wrong Answer submissions. The two tools did not differ (Hodges--Lehmann estimate $-$0.012, 95\% CI [$-$0.118, 0.074], Wilcoxon p = 0.483, N = 22). Two problems showed up in the data. Most newly written tests were invalid, mainly because of wrong expected values. With its default configuration, PITest stopped without producing a report in 58 of 96 runs, because the test suite contained tests that failed on the program under test. These preliminary results suggest educators should teach students to write valid tests and adjust the PITest configuration, instead of choosing between report types. The study is exploratory and needs confirmation with real students and the planned sample of 48.
\end{abstract}
